% English + Spanish/Catalan abstracts. UAB language requirement TBD.
\chapter*{Abstract}
\addcontentsline{toc}{chapter}{Abstract}

Randomized controlled trials represent a fast-growing and critical tool in empirical
economics, boasting a literature that is vast, still growing, and has recently birthed
its own Nobel prizes. Much of the focus on this literature, however, is on the internal
validity of the experimental design, as opposed to the external validity of results for
policy relevance. This thesis concerns itself with two key components of external
validity: \emph{representativity} and \emph{ecological validity}. It develops novel
methodologies for leveraging the particular properties of digital ad platforms to
address both of these components: (i) ecological validity when studying the impact of
digital interventions and (ii) representativity more broadly, through optimizing
participant recruitment subject to budget constraints.

Three studies develop this argument. Chapter~1 reframes representativity as a continuous
quantity to be optimized under a budget constraint and develops an adaptive stratified
recruitment algorithm that learns per-stratum costs in real time and reallocates spend
toward the strata that most reduce estimation error; validated in the United States
against the General Social Survey, Current Population Survey, and Pew benchmarks, it
achieves a mean absolute deviation of roughly 6.1 percentage points --- outperforming a
parallel Prolific sample --- at \$0.32 per respondent-question, and has since been
deployed across 33 studies in 23 countries. Chapter~2 introduces the \emph{feed
experiment}, a natural field experiment that delivers a malaria-prevention campaign as
real ads in participants' news feeds and measures behaviour through a channel unconnected
to the advertiser: a district-level cluster-randomized trial of the at-scale campaign in
India finds limited average effects but detectable impacts concentrated among
higher-socioeconomic-status and urban households, while the paired feed experiment
indicates the ad content itself can shift bed-net use across the at-risk population,
locating the binding constraint in delivery rather than persuasion. Chapter~3 evaluates
UNICEF's Bebbo parenting app through a randomized encouragement design in Serbia and
Bulgaria, finding no significant effect on parenting knowledge, attitudes, or practices
under either intent-to-treat or treatment-on-the-treated analysis; with 76\% of users
abandoning the app after the first day, the null reflects the product's failure to engage
rather than the design. Together they trace the limits and possibilities of empirical
economics conducted through digital advertising.


\cleardoublepage
